% Generated from paper/arxiv_ready.md; edit that source and rebuild.
% Options for packages loaded elsewhere
\PassOptionsToPackage{unicode}{hyperref}
\PassOptionsToPackage{hyphens}{url}
\PassOptionsToPackage{dvipsnames,svgnames,x11names}{xcolor}
\documentclass[
  11pt,
  a4paper,
]{article}
\usepackage{xcolor}
\usepackage[margin=25mm]{geometry}
\usepackage{amsmath,amssymb}
\setcounter{secnumdepth}{-\maxdimen} % remove section numbering
\usepackage{iftex}
\ifPDFTeX
  \usepackage[T1]{fontenc}
  \usepackage[utf8]{inputenc}
  \usepackage{textcomp} % provide euro and other symbols
\else % if luatex or xetex
  \usepackage{unicode-math} % this also loads fontspec
  \defaultfontfeatures{Scale=MatchLowercase}
  \defaultfontfeatures[\rmfamily]{Ligatures=TeX,Scale=1}
\fi
\usepackage{lmodern}
\ifPDFTeX\else
  % xetex/luatex font selection
\fi
% Use upquote if available, for straight quotes in verbatim environments
\IfFileExists{upquote.sty}{\usepackage{upquote}}{}
\IfFileExists{microtype.sty}{% use microtype if available
  \usepackage[]{microtype}
  \UseMicrotypeSet[protrusion]{basicmath} % disable protrusion for tt fonts
}{}
\makeatletter
\@ifundefined{KOMAClassName}{% if non-KOMA class
  \IfFileExists{parskip.sty}{%
    \usepackage{parskip}
  }{% else
    \setlength{\parindent}{0pt}
    \setlength{\parskip}{6pt plus 2pt minus 1pt}}
}{% if KOMA class
  \KOMAoptions{parskip=half}}
\makeatother
% definitions for citeproc citations
\NewDocumentCommand\citeproctext{}{}
\NewDocumentCommand\citeproc{mm}{%
  \begingroup\def\citeproctext{#2}\cite{#1}\endgroup}
\makeatletter
 % allow citations to break across lines
 \let\@cite@ofmt\@firstofone
 % avoid brackets around text for \cite:
 \def\@biblabel#1{}
 \def\@cite#1#2{{#1\if@tempswa , #2\fi}}
\makeatother
\newlength{\cslhangindent}
\setlength{\cslhangindent}{1.5em}
\newlength{\csllabelwidth}
\setlength{\csllabelwidth}{3em}
\newenvironment{CSLReferences}[2] % #1 hanging-indent, #2 entry-spacing
 {\begin{list}{}{%
  \setlength{\itemindent}{0pt}
  \setlength{\leftmargin}{0pt}
  \setlength{\parsep}{0pt}
  % turn on hanging indent if param 1 is 1
  \ifodd #1
   \setlength{\leftmargin}{\cslhangindent}
   \setlength{\itemindent}{-1\cslhangindent}
  \fi
  % set entry spacing
  \setlength{\itemsep}{#2\baselineskip}}}
 {\end{list}}
\usepackage{calc}
\newcommand{\CSLBlock}[1]{\hfill\break\parbox[t]{\linewidth}{\strut\ignorespaces#1\strut}}
\newcommand{\CSLLeftMargin}[1]{\parbox[t]{\csllabelwidth}{\strut#1\strut}}
\newcommand{\CSLRightInline}[1]{\parbox[t]{\linewidth - \csllabelwidth}{\strut#1\strut}}
\newcommand{\CSLIndent}[1]{\hspace{\cslhangindent}#1}
\setlength{\emergencystretch}{3em} % prevent overfull lines
\providecommand{\tightlist}{%
  \setlength{\itemsep}{0pt}\setlength{\parskip}{0pt}}
% Layout additions for the standalone paper generated from arxiv_ready.md.
\usepackage{microtype}
\allowdisplaybreaks[1]
\setlength{\emergencystretch}{2em}
\clubpenalty=10000
\widowpenalty=10000
\displaywidowpenalty=10000
\usepackage{bookmark}
\IfFileExists{xurl.sty}{\usepackage{xurl}}{} % add URL line breaks if available
\urlstyle{same}
% fallback for those not using the hyperref driver hyperxmp:
\makeatletter
\@ifundefined{xmpquote}{\newcommand{\xmpquote}[1]{#1}}{}
\makeatother
\hypersetup{
  pdftitle={Stability of a prescribed repeatable completion of a quantum measurement},
  pdfauthor={Christopher Gardner},
  colorlinks=true,
  linkcolor={blue},
  filecolor={Maroon},
  citecolor={blue},
  urlcolor={blue},
  pdfcreator={LaTeX via pandoc}}

\title{Stability of a prescribed repeatable completion of a quantum
measurement}
\author{Christopher Gardner}
\date{1 November 2026}

\begin{document}
\maketitle

\textbf{External-review draft.} The argument below gives an explicit
quantitative specialization of established constructions and
inequalities. The supplied reduction, its usefulness, and publication
priority remain subjects for independent review.

\section{Abstract}\label{abstract}

Let \(E=(E_j)\) be a finite positive-operator-valued measurement (POVM),
and let \(P_j\) be its full spectral sectors at or above \(1-\delta\),
where \(0<\delta<1/2\). Assume that each nonzero effect has a nonzero
sector and that the sector synthesis operator is injective, with least
singular value \(\alpha\). Its unweighted polar isometry determines
orthogonal sectors \(Q_j\). We prove that the prescribed completion
\(F_j=Q_j+RE_jR\), where \(R=I-\sum_jQ_j\), admits a repeatable
instrument and satisfies \[\|\mathcal M_E-\mathcal M_F\|_\diamond
\le 2\min\{1,\sqrt{1-\alpha^2}+\sqrt\delta\}.\] The proof reduces to
polar geometry, a classical Jordan-product inequality, and pure-state
trace distance after an explicit labelled-dilation calculation. For
orthogonal input sectors, we determine the exact worst-case envelope
\(2\sqrt\delta/(2-\sqrt\delta)\) and exhibit an attaining family. Exact
rational certificates show both a strict reference advantage within this
subclass and the failure of a natural joint extension outside it. A
coupled simplex family gives lower bounds for imperfect packing. A
corollary corrects separated output records to an exactly idempotent
channel. Both sector constructions are established least-squares
measurements; no new general inequality or literature priority is
claimed. The sharp joint bound for general packing remains open.

\section{1. Question and scope}\label{question-and-scope}

A finite POVM can admit a repeatable instrument without being a
projection-valued measurement (PVM). The relevant criterion is that
every nonzero effect have eigenvalue one (Busch et al. 1995). Such
measurements have orthogonal unit eigenspaces and may retain an
arbitrary POVM on their common orthogonal complement. We study a
particular way to select a nearby measurement in this class while
retaining the input space, labels, and compressed residual effects.

The construction uses the polar factor of a synthesis operator, a
standard orthogonalization device (Higham 1986). Its sector vectors are
exactly the unweighted least-squares measurement of Eldar--Forney; the
weighted counterpart below is their weighted version (Eldar and Forney
2001). The issue addressed here is quantitative control of a prescribed
residual extension in the full diamond norm. The hypothesis concerns
near-unit sectors and their simultaneous linear independence; it is not
a global almost-projectivity assumption.

The distinction from earlier work is limited. De la Salle's
orthogonalization proof also contains a closely related weighted polar
construction (Salle 2022). Dilation continuity is standard (Kretschmann
et al. 2008), and the noncommuting matrix inequality used below is a
special case of the Strang--de Bruijn Jordan-product bounds (Bruijn
1980). No additional general operator lemma is needed: the
measurement-specific work is the labelled dilation and residual
compression, followed by parameter substitution. Section 5 makes the
reduction explicit. This is a quantitative specialization with a
complete argument.

\section{2. Setting and main
statement}\label{setting-and-main-statement}

Let \(\mathcal H\ne\{0\}\) be a finite-dimensional complex Hilbert
space. Write \(\|\cdot\|\) for operator norm and \(\|\cdot\|_1\) for
trace norm. A POVM \(E=(E_j)_{j=1}^m\), \(m\ge1\), consists of positive
operators summing to \(I\). Its measurement channel is
\[\mathcal M_E(\rho)=\sum_j\operatorname{Tr}(E_j\rho)|j\rangle\langle j|.\]
We use the unhalved diamond norm; all state distances below likewise use
the unhalved trace norm. For a difference of channels the diamond norm
equals the supremum of output trace-norm difference over density
operators with an arbitrary reference system; reference-assisted pure
inputs suffice (Watrous 2018, Theorem 3.51). Consequently its maximum is
two. Channels are completely positive and trace preserving (CPTP); their
adjoints are unital completely positive (UCP).

Fix \(0<\delta<1/2\) and set \[P_j=1_{[1-\delta,1]}(E_j),\qquad
\mathcal K_j=\operatorname{ran}P_j,\qquad
\mathcal K=\bigoplus_j\mathcal K_j.
\tag{1}\] Assume \(P_j\ne0\) whenever \(E_j\ne0\). A zero effect has
\(P_j=0\) and may retain its label. Define
\[T:\mathcal K\longrightarrow\mathcal H,\qquad
T(x_1,\ldots,x_m)=\sum_jx_j,
\qquad \alpha=\sigma_{\min}(T)>0.
\tag{2}\] Thus \(T\) is injective. Write \(T=UH\), with \(H=|T|\), and
let \(C_j\) be the coordinate projection of \(\mathcal K\) onto
\(\mathcal K_j\). Put \[Q_j=UC_jU^*,\qquad S=UU^*,\qquad R=I-S,\qquad
F_j=Q_j+RE_jR.
\tag{3}\]

\textbf{Theorem 1 (prescribed completion).} Under (1)--(2), the \(Q_j\)
are mutually orthogonal, have the ranks of \(P_j\), and obey
\[\|P_j-Q_j\|\le a,\qquad a:=\sqrt{1-\alpha^2}.
\tag{4}\] The family \(F\) in (3) is a POVM with full unit eigenspaces
\(Q_j\), and \[\begin{gathered}
F_iQ_j=\delta_{ij}Q_j,\\
\gamma_F:=\min_j\left[1-\|(I-Q_j)F_j(I-Q_j)\|\right]\ge\delta.
\end{gathered}
\tag{5}\] It admits a repeatable instrument and satisfies
\[\boxed{\|\mathcal M_E-\mathcal M_F\|_\diamond
\le 2\min\{1,a+\sqrt\delta\}.}
\tag{6}\] More precisely, writing \(d_0=\sqrt{1-\delta}\) and
\(M=\|H\|\), the sufficient condition \[(M-\alpha)(1-d_0)\le4\alpha d_0
\tag{7}\] implies \[\|\mathcal M_E-\mathcal M_F\|_\diamond
\le 2\sqrt{1-\alpha^2(1-\delta)}.
\tag{8}\]

The constants contain no explicit dimension, outcome-count, or
sector-rank factor. The domain remains restrictive:
\(\sum_j\operatorname{rank}P_j\le\dim\mathcal H\), so there are at most
\(\dim\mathcal H\) nonzero outcomes. The estimate does not apply to
generic overcomplete POVMs.

\section{3. Proof of Theorem 1}\label{proof-of-theorem-1}

\subsection{3.1. Geometry and the residual
gap}\label{geometry-and-the-residual-gap}

The polar factor \(U\) is an isometry with
\(\operatorname{ran}U=\operatorname{ran}T=\sum_j\operatorname{ran}P_j\).
Thus \(P_jR=0\), and the \(Q_j\) are orthogonal and rank preserving. A
unit vector supported in one nonzero coordinate sector has \(\|Tx\|=1\),
proving \(\alpha\le1\). The spectral inequality \(E_j\ge(1-\delta)P_j\)
gives \[(1-\delta)TT^*=(1-\delta)\sum_jP_j\le I,
\qquad M\le d_0^{-1}.
\tag{9}\]

For a nonzero sector let \(\iota_j:\mathcal K_j\to\mathcal K\) be the
coordinate inclusion, \(T_j=T\iota_j\), and \(U_j=U\iota_j\). Both
\(T_j\) and \(U_j\) are isometries, and
\[T_j^*U_j=\iota_j^*H\iota_j\ge\alpha I_{\mathcal K_j}.\] Hence their
overlap has smallest singular value at least \(\alpha\). It follows that
\(\|(I-P_j)Q_j\|\) and \(\|(I-Q_j)P_j\|\) are at most \(a\). For two
orthogonal projections their difference norm is the maximum of these two
quantities, proving (4). One may verify this identity by squaring the
difference and taking its two diagonal blocks relative to the first
projection. Zero sectors are immediate.

Each \(F_j\) is positive, and \[\sum_jF_j=S+R\Bigl(\sum_jE_j\Bigr)R=I.\]
Equation (5)'s sector action follows from orthogonality. Fullness of the
spectral sector gives \[E_j\le P_j+(1-\delta)(I-P_j),\qquad
RE_jR\le(1-\delta)R.
\tag{10}\] On \(S\), the effect \(F_j\) is \(Q_j\); on \(R\), its norm
is at most \(1-\delta\). Thus \(Q_j\) is its full unit eigenspace and
(5) holds, including zero effects. Choosing a density operator
\(\tau_j\) supported in each nonzero \(Q_j\) gives the instrument
\(\mathcal I_j(\rho)=\operatorname{Tr}(F_j\rho)\tau_j\). For a zero
effect set \(\mathcal I_j=0\). Since
\(\operatorname{Tr}(F_i\tau_j)=\delta_{ij}\), a repetition returns the
same outcome.

\subsection{3.2. An explicit labelled
dilation}\label{an-explicit-labelled-dilation}

On \(\mathcal Y=\bigoplus_j\mathcal H\), let \(\Lambda_j\) project onto
label \(j\), and define \[\begin{aligned}
W:\mathcal H\to\mathcal Y,&\quad Wz=(\sqrt{E_j}z)_j,\\
J:\mathcal K\to\mathcal Y,&\quad J(x_j)_j=(x_j)_j.
\end{aligned}
\tag{11}\] In particular, \(J\) embeds the original sector vectors. Both
\(W\) and \(J\) are isometries. Spectral commutation allows the block
operator \[D=\bigoplus_j\sqrt{E_j}\big|_{\mathcal K_j},
\qquad d_0I\le D\le I.\] Directly from their domains,
\[W^*J=TD=UHD,\qquad
J^*\Lambda_jWR=0\quad\text{for every }j.
\tag{12}\] Indeed, the only possible nonzero coordinate of the second
expression is \(P_j\sqrt{E_j}R=\sqrt{E_j}P_jR=0\).

Set \[V=JU^*+WR.
\tag{13}\] The cross terms in \(V^*V\) vanish, so \(V^*V=S+R=I\). The
labelwise identities in (12), together with \(J^*\Lambda_jJ=C_j\), give
\[V^*\Lambda_jV=UC_jU^*+RE_jR=F_j,\qquad
W^*V=UHD\,U^*+R.
\tag{14}\] These identities concern the particular completion (3).

\subsection{3.3. The overlap estimate}\label{the-overlap-estimate}

\textbf{Lemma 2 (classical Jordan-product specialization).} If
\(\alpha I\le H\le MI\), \(d_0I\le D\le I\), and (7) holds, then
\(\operatorname{Re}(HD)\ge\alpha d_0I\).

\emph{Proof.} Set \(X=H-\alpha I\), \(Y=D-d_0I\), \(h=M-\alpha\), and
\(r=1-d_0\). For \(s>0\), positivity of \((\sqrt{s}X+s^{-1/2}Y)^2\)
implies \[XY+YX\ge-sX^2-s^{-1}Y^2\ge-shX-(r/s)Y.\] Therefore
\[\operatorname{Re}(HD)-\alpha d_0I
\ge(d_0-sh/2)X+(\alpha-r/(2s))Y.\] For \(h>0\), take \(s=2d_0/h\); both
coefficients are nonnegative by (7). For \(h=0\), \(H=\alpha I\) and the
assertion is immediate. \(\square\)

For attribution, put \(u=(M-\alpha)/(M+\alpha)\) and
\(v=(1-d_0)/(1+d_0)\). Condition (7) is exactly \(u+v\le1\); the
corresponding branch of de Bruijn's Theorem 3.2 and Sections 4--5 gives
the lower bound \((M+\alpha)(1+d_0)(1-u)(1-v)/4=\alpha d_0\) (Bruijn
1980). The proof above is included for completeness.

By (14) and \(\alpha d_0\le1\), Lemma 2 yields
\[\operatorname{Re}(W^*V)\ge\alpha d_0S+R\ge\alpha d_0I.\] For any
normalized \(\psi\in\mathcal H\otimes\mathcal A\), the dilated vectors
\((W\otimes I_{\mathcal A})\psi\) and \((V\otimes I_{\mathcal A})\psi\)
have overlap of modulus at least \(\alpha d_0\). The trace distance of
their pure density operators is consequently at most
\(2\sqrt{1-\alpha^2d_0^2}\). Apply to both the same CPTP label readout
\(\mathcal L(X)=\sum_j\operatorname{Tr}(\Lambda_jX)|j\rangle\langle j|\),
with the identity on \(\mathcal A\). Trace-norm contractivity and the
reference-assisted characterization prove (8).

To obtain (6), write \(b=\sqrt\delta\). If \(a+b\ge1\), use the
universal bound two. Otherwise \(a<1-b\), so \(\alpha^2>2b-b^2>b^2\),
hence \(\alpha>b\). By (9), \[(M-\alpha)(1-d_0)\le M(1-d_0)
\le\frac{\delta}{d_0(1+d_0)}
\le\delta<4\alpha d_0.\] Here \(d_0>1/\sqrt2\), and \(b<4d_0\). Thus (7)
holds, and \[2\sqrt{1-\alpha^2(1-\delta)}
=2\sqrt{a^2+\alpha^2\delta}
\le2(a+b).\] This completes the proof of Theorem 1. \(\square\)

\section{4. Correction of supplied
records}\label{correction-of-supplied-records}

Let \(\sigma_j\) be density operators on the same \(\mathcal H\), and
set \[\Phi(\rho)=\sum_j\operatorname{Tr}(E_j\rho)\sigma_j,\qquad
\eta=\|\Phi^2-\Phi\|_\diamond.\] For this section discard zero effects
and their unused records. Assume at least two remaining labels and
define \[\begin{gathered}
\kappa=\min_j\frac12\inf_{\omega\in\operatorname{conv}\{\sigma_i:i\ne j\}}
\|\sigma_j-\omega\|_1>0,\\
e=\max_j(1-\lambda_{\max}(E_j)).
\end{gathered}
\tag{15}\] The main hypotheses give \(0\le e\le\delta<1/2\). With one
remaining label the channel is constant and already idempotent.

\textbf{Corollary 3 (record correction).} Under Theorem 1 and (15),
there are explicit states \(\tau_j\) supported in \(Q_j\) such that
\(\Pi(\rho)=\sum_j\operatorname{Tr}(F_j\rho)\tau_j\) is exactly
idempotent. For any valid measurement movement bound \(B_{\rm move}\),
\[\boxed{\|\Phi-\Pi\|_\diamond\le
\min\left\{2,\ B_{\rm move}+2a+
\sqrt{\frac{2\eta}{\kappa\delta(1-2e)}}\right\}.}
\tag{16}\] In particular, \[\|\Phi-\Pi\|_\diamond\le
\min\left\{2,\ 4a+2\sqrt\delta+
\sqrt{\frac{2\eta}{\kappa\delta(1-2\delta)}}\right\}.
\tag{17}\]

\emph{Proof.} Put \(v_j=\Phi(\sigma_j)-\sigma_j\) and
\(D_{\rm rec}=\max_j\|v_j\|_1\). The action of \(\Phi^2-\Phi\) on a
reference-assisted state is \(\sum_jv_j\otimes B_j\), where the
conditional operators \(B_j\) are positive and
\(\sum_j\operatorname{Tr}B_j=1\). Thus \(\eta\le D_{\rm rec}\). Choose a
maximizing label and an input whose probability of that label is at
least \(1-e\). The reverse triangle inequality gives
\[(1-2e)D_{\rm rec}\le\eta\le D_{\rm rec}.
\tag{18}\] Writing \(a_j=\operatorname{Tr}(E_j\sigma_j)\), when
\(a_j<1\) the other outcomes yield a state \(\omega_j\) in the competing
convex hull with \(v_j=(1-a_j)(\omega_j-\sigma_j)\). The case \(a_j=1\)
is immediate. Hence \[2\kappa(1-a_j)\le\|v_j\|_1,\qquad
1-\operatorname{Tr}(P_j\sigma_j)\le
q:=\frac{\eta}{2\kappa\delta(1-2e)},
\tag{19}\] where the second inequality uses (10).

For a purification of \(\sigma_j\), the triangle inequality and (4)
imply \[\sqrt{1-\operatorname{Tr}(Q_j\sigma_j)}
\le\sqrt{1-\operatorname{Tr}(P_j\sigma_j)}+\|P_j-Q_j\|
\le\sqrt q+a.\] If \(z_j=\operatorname{Tr}(Q_j\sigma_j)>0\), set
\(\tau_j=Q_j\sigma_jQ_j/z_j\). Projecting the purification gives a
normalized vector with overlap \(\sqrt{z_j}\) with the original; the
pure-state trace distance formula and partial trace therefore give
\[\|\sigma_j-\tau_j\|_1\le2\sqrt{1-z_j}\le2\sqrt q+2a.
\tag{20}\] This is the standard normalized gentle-projection estimate
(Watrous 2018, Corollary 3.15 and Theorem 3.33). If \(z_j=0\), the
preceding inequality forces \(\sqrt q+a\ge1\); choose
\(\tau_j=Q_j/\operatorname{rank}Q_j\), which satisfies (20).

Now \(\operatorname{Tr}(F_i\tau_j)=\delta_{ij}\), so \(\Pi^2=\Pi\).
Keeping the records \(\sigma_j\) while changing \(E\) to \(F\) costs at
most \(B_{\rm move}\), by common CPTP post-processing. Keeping \(F\)
while changing its records costs at most
\(\max_j\|\sigma_j-\tau_j\|_1\), also for inputs with a reference.
Combining these facts proves (16). Theorem 1 and \(e\le\delta\) give
(17). \(\square\)

For a concrete nontrivial regime, take \(0<\varepsilon<1/16\) and
\(E_0=\operatorname{diag}(1-\varepsilon,\varepsilon)\), \(E_1=I-E_0\),
and the two computational-basis records. Then \(\alpha=\kappa=1\),
\(e=\varepsilon\), \(\eta=2\varepsilon(1-2\varepsilon)\), and the
constructed channel is complete dephasing, at distance \(2\varepsilon\).
Choosing \(\delta=\sqrt\varepsilon\) in (16) with
\(B_{\rm move}=2\sqrt\delta\) gives \(4\varepsilon^{1/4}\), which is
below two for \(\varepsilon<1/16\). The certificate is useful but can be
very conservative.

\section{5. Relationship to existing
constructions}\label{relationship-to-existing-constructions}

Choose an eigenbasis within every full sector and write the concatenated
unit vectors as the columns of \(T\). In these coordinates
\(D=\operatorname{diag}(\sqrt{\lambda_a})\), with \(\lambda_a\) the
selected effect eigenvalues. Eldar--Forney's unweighted and weighted
least-squares formulas are exactly \[U=T(T^*T)^{-1/2},\qquad
\widehat U=TD(DT^*TD)^{-1/2}\] (Eldar and Forney 2001, Theorem 2 and
Section V). The weighted formula is also the square-root measurement for
priors proportional to \(\lambda_a\). Grouping columns by sector
recovers \(Q_j\) and \(\widehat Q_j\), including higher-rank sectors.
These formulas optimize vector least-squares objectives; they do not
assert optimal distance from the original measurement channel. The
extension \(RE_jR\) and its channel comparison must still be checked.

\textbf{Proposition 4 (weighted comparison).} Keep the hypotheses and
notation of Theorem 1, but let \(TD=\widehat U|TD|\). Then, without
condition (7), \[\begin{gathered}
\widehat F_j=\widehat UC_j\widehat U^*+RE_jR,\\
\|\mathcal M_E-\mathcal M_{\widehat F}\|_\diamond
\le2\sqrt{1-\alpha^2(1-\delta)}.
\end{gathered}
\tag{21}\] \emph{Proof.} The invertibility of \(D\) gives
\(\operatorname{ran}(TD)=\operatorname{ran}T\) and
\(|TD|\ge\alpha d_0I\). Thus \(\widehat U\) has the same range
projection \(S\). Equations (12)--(14), with
\(\widehat V=J\widehat U^*+WR\), give the isometry and effect
identities, while
\[W^*\widehat V=\widehat U|TD|\widehat U^*+R\ge\alpha d_0I.\] The same
reference-assisted overlap argument proves (21). \(\square\)

This is a reconstructed application of established polar alignment, not
a claim that a cited paper states (21) verbatim. Its connection to de la
Salle's proof is explicit. On \(S\mathcal H\), put \(A_j=SE_jS\). The
\(P_j\) commute with their \(A_j\), have total rank
\(\dim S\mathcal H\), and satisfy
\[\sum_jP_jA_j\ge(1-\delta)\alpha^2I_{S\mathcal H}.\] They can therefore
replace the state-selected projections in the finite proof's polar step.
That column becomes \(x=(P_j\sqrt{A_j})_j=DT^*\); its adjoint is \(TD\)
(Salle 2022, Lemma 3.1 and Section 3). The residual extension and
overlap calculation above make the application uniform. Broad existence
of a nearby same-label repeatable measurement consequently has a short
derivation from this machinery. The weighted construction generally
differs from (3); (21) does not prove Theorem 1's estimates for that
prescribed map.

KSW's continuity theorem supplies the generic dilation-to-channel
comparison (Kretschmann et al. 2008, Proposition 3). In tracking the
sharper constant in (8), our proof uses pure-state trace distance
directly: overlap at least \(c\) gives \(2\sqrt{1-c^2}\); applying only
\(2\|W-V\|\) would give \(2\sqrt{2(1-c)}\). Both are established
estimates. Kotowski--Oszmaniec's dimension-deficient Naimark
construction proves projective simulability for a different
complement-noise completion, without requiring unit eigenspaces or
asserting the present diamond estimate (Kotowski and Oszmaniec 2025,
Theorem 6).

The exact record endpoint also has established structure. On the algebra
of active labels define \(\Delta(f)=\sum_jf_jF_j\) and
\(\Upsilon(X)=(\operatorname{Tr}(\tau_jX))_j\). They are UCP maps with
\(\Upsilon\Delta=I\) and \(\Pi^*=\Delta\Upsilon\), a commutative
instance of exact encoding/decoding structure (Kitaev 2025, Theorem
1.2). Kitaev's Theorem 12.3 gives general approximate factors; it does
not directly supply this particular exact correction of the original
effects and records. The additional specification here is quantitative
and applies to a substantially narrower class.

The reduction of the unweighted estimate is complete: (12)--(14)
identify its dilation, the classical matrix bound supplies the overlap,
and the pure-state formula and scalar case split prove (6). The sector
estimate is principal-angle geometry. The record proof adds the reverse
triangle inequality, convex-hull separation, and normalized gentle
projection. Thus the bounds require no additional general lemma. The
remaining questions are whether this explicit specialization is already
documented and whether its quantitative control and limitations are
useful. The certified comparison below distinguishes the two targets,
but does not establish originality.

\section{6. Counterexamples and
limitations}\label{counterexamples-and-limitations}

\subsection{6.1. Repeatable approximation differs from projective
approximation}\label{repeatable-approximation-differs-from-projective-approximation}

On \(\mathbb C^3\), let
\[E_0(t)=\operatorname{diag}(1-t,t,1/2),\quad E_1(t)=I-E_0(t),
\qquad0<t<1/2.\] For \(\delta=t\), the sectors are the first two
coordinate lines and \(\alpha=1\). The prescribed completion is
\(F_0=\operatorname{diag}(1,0,1/2)\), \(F_1=I-F_0\), with distance
\(2t\). However, its input measurement has distance exactly one to the
nearest binary PVM. For any projection \(L\), ordered-eigenvalue
perturbation gives \(\|E_0(t)-L\|\ge1/2\), because the middle eigenvalue
of \(E_0(t)\) is \(1/2\) whereas that of \(L\) is zero or one. Binary
measurement-channel distance is \(2\|E_0(t)-L\|\); the upper bound one
is attained by \(L=\operatorname{diag}(1,0,0)\). The binary norm
identity follows from the norm of the Hermitian functional
\(\rho\mapsto\operatorname{Tr}((E_0-L)\rho)\); an extremal eigenvector
attains it, and reference systems do not increase it. Also
\(\|I-\sum_jE_j(t)^2\|=1/2\).

Thus small \(\delta\) and perfect packing do not imply global
near-projectivity. The example does not preclude the packed-subspace
reduction in Section 5.

\subsection{6.2. A square-root threshold
obstruction}\label{a-square-root-threshold-obstruction}

\textbf{Proposition 5.} For every \(0<\delta<1/10\), there is a
three-outcome POVM on \(\mathbb C^4\) satisfying Theorem 1 with
\(\alpha=1\), such that every repeatable POVM \(G\) on the same space
and with the same three labels obeys
\[\|\mathcal M_E-\mathcal M_G\|_\diamond\ge
\frac{\sqrt\delta}{4\sqrt3}.
\tag{22}\]

\emph{Proof.} In the basis \(e_1,\ldots,e_4\), put \(t=\sqrt\delta/4\)
and \[\begin{gathered}
E_1=
\begin{pmatrix}
1-\delta/2&0&0&0\\
0&\delta/2&0&t\\
0&0&0&0\\
0&t&0&1/4
\end{pmatrix},\\
E_2=\operatorname{diag}(\delta/2,1-\delta,\delta/2,1/2),
\quad E_3=I-E_1-E_2.
\end{gathered}
\tag{23}\] The nontrivial low \(2\times2\) blocks of \(E_1,E_3\) have
determinant \(\delta/16>0\) and trace \(1/4+\delta/2<3/10\). The
remaining entries are nonnegative. Thus these are POVM effects, and
their full near-unit sectors are respectively
\(\mathbb Ce_1,\mathbb Ce_2,\mathbb Ce_3\); hence \(\alpha=1\).

Set \(\rho_0=\max_j\|E_j-G_j\|\). If \(\rho_0\le1/4\), then \(G_2\ne0\),
since \(\|E_2\|=1-\delta>9/10\). Repeatability supplies a unit vector
\(v\) with \(G_2v=v\), and positivity of the other effects gives
\(G_1v=0\). Write \(v=a_2e_2+b_2\), \(b_2\perp e_2\). Since
\(I-E_2\ge I/2\) on \(e_2^\perp\), \[\|(I-E_2)v\|\le\rho_0
\quad\Longrightarrow\quad
\|b_2\|\le2\rho_0,\quad |a_2|\ge\sqrt3/2.\] Moreover,
\(\|E_1v\|\le\rho_0\), \(\|E_1\|\le1\), and \(\|E_1e_2\|\ge t\).
Therefore \[|a_2|t\le\|a_2E_1e_2\|
\le\rho_0+\|b_2\|\le3\rho_0,\] so \(\rho_0\ge t/(2\sqrt3)\). For any two
POVMs, a state maximizing one effect difference gives
\(\|\mathcal M_E-\mathcal M_G\|_\diamond\ge2\rho_0\): the other
probability differences sum to its negative. This proves (22). If
\(\rho_0>1/4\), the same conclusion follows from distance greater than
\(1/2\). \(\square\)

The proposition establishes the necessity of square-root threshold order
for this same-space, same-label target class. It does not establish
optimality of the constant two, the dependence on \(\alpha\), or the
joint bound.

\subsection{6.3. Neither completion uniformly dominates the
other}\label{neither-completion-uniformly-dominates-the-other}

\textbf{Proposition 6 (computer-assisted comparison).} On
\(\mathbb C^4\), set \[\begin{gathered}
E_1(t)=\frac1{40}\begin{pmatrix}
32&0&3&1\\0&2&2&2\\3&2&4&4\\1&2&4&8
\end{pmatrix},\qquad
E_2(t)=\frac1{40}\begin{pmatrix}
4&2&-2&t\\2&34&-3&-1\\-2&-3&2&-2\\t&-1&-2&12
\end{pmatrix},\\
E_3(t)=I-E_1(t)-E_2(t),\qquad 2\le t\le4,\qquad\delta=1/4.
\end{gathered}
\tag{24}\] This family satisfies Theorem 1 throughout the interval, with
sector ranks \((1,1,1)\) and \(\alpha(t)>9/10\). Let \(\nu_u(t)\) and
\(\nu_w(t)\) be the full diamond distances from the original measurement
to the unweighted and weighted completions, and put
\(g(t)=\nu_w(t)-\nu_u(t)\). Exact certificates give
\[g(2)>5\times10^{-4},\qquad g(4)<-9\times10^{-5}.
\tag{25}\] The distances are continuous, and at least one equality point
satisfies \[2.33690021<t_*<2.33690023,\qquad g(t_*)=0.
\tag{26}\] Uniqueness and absence of other equality points are not
asserted.

\emph{Certificate argument.} The accompanying rational data verify
positivity at \(t=2,4\), hence throughout the affine family. Designated
diagonal entries give lower bounds \(4/5,17/20,17/20\) for the top
eigenvalues; the sums of the remaining eigenvalues are at most
\(7/20,9/20,7/10\), respectively. These bounds isolate exactly one
eigenvalue above \(3/4\) in each effect. The uniform spectral-residual
and polar certificate proves \(\alpha(t)>9/10\) on the entire interval,
not merely on a finite grid.

For each certified parameter and completion, the data specify rational
effects \(G_j\), a rigorous bound \(\epsilon\ge\sum_j\|F_j-G_j\|\), and
exact primal/dual feasible SDP witnesses with values \(L,U\) for
\(\mathcal M_E-\mathcal M_G\). Thus \(L-\epsilon\le\nu\le U+\epsilon\),
including reference systems. These enclosures prove (25) and strict
opposite signs at the two endpoints in (26). Uniform spectral isolation
and synthesis injectivity imply continuity, so the intermediate value
theorem proves (26). The
\href{https://github.com/christophergardner-star/RC-Discovery/blob/d706872d89503f5c6f9e59d8f91b518bec9991b1/research/COMPLETION_CROSSOVER_2026-11-01.md}{certificate
supplement} links the data, the analytic error bounds and uniform
argument, and the exact standard-library verifier. It is part of the
computer-assisted proof of this proposition. \(\square\)

The additive theorem bound stays below \(1+\sqrt{19}/5<2\) along this
path. A stronger weighted upper envelope therefore does not imply a
smaller actual weighted distance. These comparisons do not prove either
completion optimal among all repeatable targets, sharpness of the upper
bounds, or literature priority of the examples.

\subsection{6.4. Other limits}\label{other-limits}

The strict accessibility condition in (18) is essential. For
\(E_0=E_1=I/2\) and orthogonal pure records, \(e=1/2\), \(\Phi\) is
constant, \(\eta=0\), but \(D_{\rm rec}=1\). This example violates the
main theorem's near-unit hypothesis; it concerns the accessibility
lemma's endpoint.

The prescribed spectral selection is not globally continuous. Fix
\(\delta=1/5\) and
\[E_0(\varepsilon)=\operatorname{diag}(1,4/5-\varepsilon,0,1/2),
\qquad E_1(\varepsilon)=I-E_0(\varepsilon).\] At \(\varepsilon=0\),
\(F_0=\operatorname{diag}(1,1,0,1/2)\). For \(0<\varepsilon<1/10\),
\(F_0(\varepsilon)=E_0(\varepsilon)\). All examples satisfy the
hypotheses with \(\alpha=1\). Input measurement distance
\(2\varepsilon\) tends to zero, while the completed measurement distance
tends to \(2/5\). A perturbation-continuity claim therefore requires
additional threshold isolation.

Finally, (6) is informative only when \(a+\sqrt\delta<1\). The record
certificate can become trivial for small \(\kappa\) or unfavorable
\(\eta/\delta\). No experimental estimation procedure for all these
parameters, nor practical advantage of the unweighted over the weighted
completion, is established here.

\section{7. Sharpness and reference
systems}\label{sharpness-and-reference-systems}

\subsection{7.1. Exact envelope for orthogonal
sectors}\label{exact-envelope-for-orthogonal-sectors}

\textbf{Theorem 7 (perfect packing).} Under Theorem 1 with \(\alpha=1\),
\[\nu:=\|\mathcal M_E-\mathcal M_F\|_\diamond
\le B(\delta):=\frac{2\sqrt\delta}{2-\sqrt\delta}.
\tag{27}\] For every \(0<\delta<1/2\), equality is attained by a
five-outcome POVM on \(\mathbb C^6\), using a system-only input. Thus
(27) is the exact worst-case envelope in this subclass, over arbitrary
finite dimensions and outcome counts.

\emph{Proof.} The diagonal blocks of \(T^*T\) are identities, so its
trace is \(\dim\mathcal K\). Its least eigenvalue is one; hence
\(T^*T=I\). The \(P_j\) are orthogonal and \(Q_j=P_j\). Write
\(S=\sum_jP_j\), \(R=I-S\), \(A_j=P_jE_jP_j\), \(\Delta_j=P_j-A_j\), and
\(\Delta=\sum_j\Delta_j\). Spectral commutation gives
\(B_j:=E_j-A_j\ge0\), with \(\sum_jB_j=L:=R+\Delta\) and
\(0\le\Delta\le\delta S\). Normalize the \(B_j\) on the support of
\(L\): \(N_j=L^{-1/2}B_jL^{-1/2}\). Allocate the kernel projection to
any label to obtain a POVM \(N\). Since \(B_j\le L\), the factorization
\(B_j=\sqrt L N_j\sqrt L\) is exact, including on the kernel. Therefore
\[\mathcal M_E-\mathcal M_F
=\mathcal M_N\bigl(\sqrt L\,\cdot\,\sqrt L-R\,\cdot\,R\bigr)
-\mathcal M_\Delta.
\tag{28}\] Here \(\mathcal M_\Delta\) is completely positive, though
generally not trace preserving.

For any unit vector \(\psi\), including a reference system, put
\(p=\|(S\otimes I)\psi\|^2\), \(q=1-p\), and
\(r=\langle\psi,(\Delta\otimes I)\psi\rangle\le\delta p\). The vectors
\(x=(\sqrt L\otimes I)\psi\) and \(y=(R\otimes I)\psi\) have squared
norms \(q+r,q\) and overlap \(q\). Their rank-two difference has trace
norm \(\sqrt{r^2+4rq}\), by its trace \(r\) and eigenvalue product
\(-rq\). This standard projected-state calculation also appears
explicitly in (Regula et al. 2026, Appendix C). Contractivity and the
trace \(r\) of the positive defect output give
\[\|[(\mathcal M_E-\mathcal M_F)\otimes\mathrm{id}](\psi\psi^*)\|_1
\le r+\sqrt{r^2+4rq}
\le\delta p+\sqrt{\delta^2p^2+4\delta p(1-p)}.
\tag{29}\] Set \(b=\sqrt\delta\), \(z=2-b\), \(B=2b/z\). Since
\(B-b^2p>0\), the identity \[(B-b^2p)^2-\bigl(b^4p^2+4b^2p(1-p)\bigr)
=4b^2(p-1/z)^2\] proves (27), with equality in the scalar optimization
at \(p=1/z\). Reference-assisted pure inputs suffice for the diamond
norm.

For attainment, put \(h=\sqrt{1-b}\), so \(z=1+h^2\), and use basis
\(e_0,e_1,e_2,e_3,e_4,e_r\). Let \(P_j=e_je_j^*\). On
\(\operatorname{span}\{e_0,e_r\}\), define \[\begin{gathered}
L_+=\frac1z\begin{pmatrix}b^2&bh\\bh&h^2\end{pmatrix},\qquad
L_-=\frac1z\begin{pmatrix}b^2h^2&-bh\\-bh&1\end{pmatrix},\\
E_0=(1-b^2)P_0,\qquad
E_1=P_1+L_+/2,\quad E_2=P_2+L_+/2,\\
E_3=P_3+L_-/2,\quad E_4=P_4+L_-/2.
\end{gathered}
\tag{30}\] The positive rank-one operators \(L_\pm\) sum to
\(b^2P_0+R\le I\). Thus the effects sum to \(I\), and every low part has
norm at most \(1/2<1-\delta\). Their full threshold sectors are exactly
the \(P_j\), including the closed-threshold eigenvalue of \(E_0\). For
\(\psi=(e_0+he_r)/\sqrt z\), the five probability differences are
\(-b^2/z,b/(2z),b/(2z),-bh^2/(2z),-bh^2/(2z)\). Their absolute sum is
\(2b/z\), proving attainment. \(\square\)

Consequently \(B(\delta)=\sqrt\delta+\delta/2+O(\delta^{3/2})\). The
optimal asymptotic coefficient is one; the best uniform coefficient of
\(\sqrt\delta\) for \(0<\delta<1/2\) is \((8+2\sqrt2)/7\), approached as
\(\delta\uparrow1/2\). This is sharpness for the prescribed completion,
not for the best freely chosen repeatable target. If \(R=0\), (29)
instead gives \(\nu\le2\delta\). The
\href{https://github.com/christophergardner-star/RC-Discovery/blob/327986e318cbfeb8632f61dbd9bad2f0dc18b3c3/theory/perfect_packing_sharpness.md}{extended
proof and exact attaining certificates} provide additional checks
independent of numerical solver tolerances.

\subsection{7.2. A reference can still improve
discrimination}\label{a-reference-can-still-improve-discrimination}

Define
\(\nu_{\rm sys}=\sup_\rho\sum_j|\operatorname{Tr}((E_j-F_j)\rho)|\).
Theorem 7's attaining family does not imply \(\nu=\nu_{\rm sys}\) for
every measurement with \(\alpha=1\). Reference-assisted measurement
discrimination is an established subject (Puchała et al. 2018); the
point here is its occurrence within this particular completion class.

\textbf{Proposition 8 (computer-assisted reference advantage).} There is
a rational eight-outcome POVM on \(\mathbb C^{10}\), satisfying Theorem
1 with \(\alpha=1\) and \(\delta=17/100\), for which
\[\nu_{\rm sys}\le0.1161<0.1208\le\nu.
\tag{31}\]

\emph{Certificate proof.} The
\href{https://github.com/christophergardner-star/RC-Discovery/blob/56f50f49db301a0f69eaed8ab5a650591a6a9aae/theory/reference_advantage.md}{specified
effects and certificate supplement} form part of the statement and
proof. With \(D_j=E_j-F_j\),
\(\nu_{\rm sys}=\max_{s\in\{-1,1\}^8}\lambda_{\max}(\sum_js_jD_j)\).
Exact positivity of all 256 matrices \(0.1161I-\sum_js_jD_j^T\) gives
the upper bound. Rational Hermitian SDP witnesses satisfy
\(\operatorname{Tr}\rho=1\), \(\rho\ge0\), \(-\rho\le X_j\le\rho\), and
\(\sum_j\operatorname{Tr}(D_j^TX_j)\ge0.1208\), proving the lower bound
by standard diamond-norm duality (Watrous 2012). The verifier also
checks every full spectral sector and the exact prescribed target. The
advantage is at least \(0.0047\); neither distance is asserted to be
exactly equal to the displayed endpoint. \(\square\)

\subsection{7.3. Imperfect packing and failed
extensions}\label{imperfect-packing-and-failed-extensions}

\textbf{Proposition 9 (coupled simplex lower bound).} For every even
\(m\ge4\) and \(0<\alpha\le1\), an \(m\)-outcome POVM on
\(\mathbb C^{m+1}\) satisfies Theorem 1 with the stated \(\alpha\),
\(\delta=1/m\), and \[\nu\ge L_m(\alpha):=
\sqrt{\frac{(m-2)[m-(m-1)\alpha^2]}{m(m-2)+m-(m-1)\alpha^2}}.
\tag{32}\] For fixed \(\alpha<1\), this lower bound tends to
\(\sqrt{1-\alpha^2}\) as \(m\to\infty\). If
\(a_m^2=1-\alpha_m^2=O(1/m)\), then
\(L_m(\alpha_m)/\sqrt{a_m^2+1/m}\to1\).

\emph{Construction argument.} The
\href{https://github.com/christophergardner-star/RC-Discovery/blob/4bb04ff6e2a5f9250a81a0093cfbb1016cbbe82d/theory/imperfect_packing_lower_bound.md}{analytic
supplement} specifies the matrices and forms part of this proof. The
synthesis is \(cI+(\alpha-c)uu^*\), with \(u=m^{-1/2}(1,\ldots,1)\) and
\(c=\sqrt{(m-\alpha^2)/(m-1)}\). Residual coupling vectors sum to zero
and are orthogonal to their respective selected eigenvectors. Their size
is fixed by a Schur complement, proving positivity. Each selected
eigenvalue is \(1-1/m\); its positive complementary block has trace
\(2/m<1-1/m\), so the full sectors are isolated. The positive synthesis
makes the polar sectors coordinate projections. Input \((u+e_r)/\sqrt2\)
has output distance exactly \(L_m(\alpha)\). The limits follow from
(32). \(\square\)

The joint sharp upper bound remains open. Two tempting extensions of
(27) are already false:
\[\nu\not\le 2(1-\alpha)+B(\delta)\quad\hbox{in general},\qquad
\nu\not\le\sqrt{1-\alpha^2+B(\delta)^2}\quad\hbox{in general}.
\tag{33}\] For the first, (32) with \(m=64\), \(\alpha=9/10\) has
squared lower bound \(80414/398097>1/9\), while the proposed upper bound
is \(1/3\). For the second, a rational seven-dimensional variant has
\(\alpha=13/27\), \(\delta=1/6\), system-only distance at least
\(1.0405\), and proposed bound at most \(1.0156\). The same supplement
links its exact effects, rational input and sign witness, and verifier.
These are counterexamples to proposed extensions, not to Theorem 1.

\section{8. Reproducibility, methodology, and review
status}\label{reproducibility-methodology-and-review-status}

The
\href{https://github.com/christophergardner-star/RC-Discovery}{RC-Discovery
repository} contains the frozen statement, analytic source notes,
numerical identity checks, reference-assisted primal/dual SDPs, negative
controls, and recorded seeds and source hashes. Its theorem-level
comparison was checked on 104 weighted constructions and six primal/dual
SDPs, alongside exact rational examples. The analytic proofs of Theorem
1, Corollary 3, and Propositions 4--5 do not rely on solver output.
Proposition 6 is computer-assisted: rational certificates and the stated
analytic error bounds establish the inequalities. Its combined audit
passes 139 independent integer-determinant positivity checks and 12
rejection controls. These checks do not prove uniqueness. Theorem 7 is
analytic and has seven additional exact attaining certificates.
Proposition 8's audit checks 289 positive-definite matrices with 5,764
integer leading minors and 13 rejection controls. The rational veto
following Proposition 9 passes 42 integer leading-minor checks and ten
rejection controls. Its universal lower-bound family is established by
the linked analytic construction, not extrapolation from numerical
tests. A separate imperfect-packing search checked 56 structured cases,
711 accepted adaptive cases, and seven SDPs; these are numerical
evidence only. The Lean development checks 13 constituent declarations,
not the full spectral and channel theorem.

This work was developed using \textbf{VIBE MATH}, an AI-assisted
workflow of conjecture generation, derivation, adversarial checking,
literature comparison, and computational support directed by Christopher
Gardner. AI systems are research tools, not authors or endorsers.
Agreement among internal AI-assisted checks is not external
verification.

The requested external assessment is concrete: check the supplied
reduction of Theorem 1 and Corollary 3 to established tools, reconstruct
the sharp subclass theorem, lower-bound family and exact comparisons,
and identify any earlier statement or application that subsumes this
quantitative specialization. The reduction requires no new general
lemma; publication priority and usefulness remain unverified.
Correctness, priority, optimality, and usefulness are separate
questions. This draft has not undergone independent specialist review
and is not represented as submission-ready.

\section*{References}\label{references}
\addcontentsline{toc}{section}{References}

\protect\phantomsection\label{refs}
\begin{CSLReferences}{1}{1}
\bibitem[\citeproctext]{ref-debruijn1980}
Bruijn, N. G. de. 1980. {``Remarks on {Hermitian} Matrices.''}
\emph{Linear Algebra and Its Applications} 32: 201--8.
\url{https://doi.org/10.1016/0024-3795(80)90017-8}.

\bibitem[\citeproctext]{ref-busch1995}
Busch, Paul, Marian Grabowski, and Pekka J. Lahti. 1995. {``Repeatable
Measurements in Quantum Theory: Their Role and Feasibility.''}
\emph{Foundations of Physics} 25 (9): 1239--66.
\url{https://doi.org/10.1007/BF02055331}.

\bibitem[\citeproctext]{ref-eldar2001}
Eldar, Yonina C., and G. David Forney Jr. 2001. {``On Quantum Detection
and the Square-Root Measurement.''} \emph{IEEE Transactions on
Information Theory} 47 (3): 858--72.
\url{https://doi.org/10.1109/18.915636}.

\bibitem[\citeproctext]{ref-higham1986}
Higham, Nicholas J. 1986. {``Computing the Polar Decomposition---with
Applications.''} \emph{SIAM Journal on Scientific and Statistical
Computing} 7 (4): 1160--74. \url{https://doi.org/10.1137/0907079}.

\bibitem[\citeproctext]{ref-kitaev2025}
Kitaev, Alexei. 2025. \emph{Almost-Idempotent Quantum Channels and
Approximate {\(C^*\)}-Algebras}. arXiv:2405.02434v2.
\url{https://doi.org/10.48550/arXiv.2405.02434}.

\bibitem[\citeproctext]{ref-kotowski2025}
Kotowski, Michał, and Michał Oszmaniec. 2025. \emph{Pretty-Good
Simulation of All Quantum Measurements by Projective Measurements}.
arXiv:2501.09339v1. \url{https://doi.org/10.48550/arXiv.2501.09339}.

\bibitem[\citeproctext]{ref-ksw2008}
Kretschmann, Dennis, Dirk Schlingemann, and Reinhard F. Werner. 2008.
{``A Continuity Theorem for {Stinespring}'s Dilation.''} \emph{Journal
of Functional Analysis} 255 (8): 1889--904.
\url{https://doi.org/10.1016/j.jfa.2008.07.023}.

\bibitem[\citeproctext]{ref-puchala2018}
Puchała, Zbigniew, Łukasz Pawela, Aleksandra Krawiec, and Ryszard
Kukulski. 2018. {``Strategies for Optimal Single-Shot Discrimination of
Quantum Measurements.''} \emph{Physical Review A} 98: 042103.
\url{https://doi.org/10.1103/PhysRevA.98.042103}.

\bibitem[\citeproctext]{ref-regula2026}
Regula, Bartosz, Ludovico Lami, and Nilanjana Datta. 2026. \emph{Tight
Relations and Equivalences between Smooth Relative Entropies}.
arXiv:2501.12447v5. \url{https://arxiv.org/abs/2501.12447v5}.

\bibitem[\citeproctext]{ref-delasalle2022}
Salle, Mikael de la. 2022. {``Orthogonalization of Positive Operator
Valued Measures.''} \emph{Comptes Rendus. Mathématique} 360: 549--60.
\url{https://doi.org/10.5802/crmath.326}.

\bibitem[\citeproctext]{ref-watrous2012}
Watrous, John. 2012. \emph{Simpler Semidefinite Programs for Completely
Bounded Norms}. arXiv:1207.5726v2.
\url{https://arxiv.org/abs/1207.5726v2}.

\bibitem[\citeproctext]{ref-watrous2018}
Watrous, John. 2018. \emph{The Theory of Quantum Information}. Cambridge
University Press. \url{https://doi.org/10.1017/9781316848142}.

\end{CSLReferences}

\end{document}
